\documentclass[letterpaper,11pt]{article}

\usepackage[letterpaper, left=0.8in, right=0.8in, top=0.7in, bottom=0.7in]{geometry}
\usepackage[hidelinks]{hyperref}
\usepackage{parskip}
\usepackage{setspace}

\setstretch{1.04}
\pagestyle{empty}

\begin{document}

\begin{center}
    {\Large \textbf{Aryan Miriyala}}\\
    \vspace{2pt}
    +1 419-315-0444 \textbar{} \href{mailto:aryanmiriyala@gmail.com}{aryanmiriyala@gmail.com} \textbar{} \href{https://www.linkedin.com/in/aryan-miriyala-7788a21b6/}{linkedin.com/in/aryan-miriyala} \textbar{} \href{https://github.com/aryanmiriyala}{github.com/aryanmiriyala}
\end{center}

\vspace{10pt}

Dear RTI Hiring Team,

I am applying for the Early-Career Software Engineer - AI Applications role because RTI's work sits where I want my software engineering to be useful: turning research questions and technical ideas into systems that people can trust. The role's focus on Python, LLMs, RAG, REST APIs, cloud services, and practical AI evaluation matches the kind of work I have been building across internships, research, and projects.

At SmartSolve, I use Codex and Claude Code for implementation planning, debugging, code review, and self-review while keeping the final judgment grounded in my own validation. I also built secure internal software with Next.js, TypeScript, PostgreSQL, Drizzle ORM, SSO, auth middleware, Docker, and Git, and scoped an AI-enabled QMS dependency register by modeling relationships among quality-management documents. That experience maps closely to RTI's need for an early-career engineer who can learn quickly, write maintainable software, work with senior engineers, and translate ambiguous requirements into working systems.

My project work adds the AI applications layer. I built FalconGraph Search, a source-grounded RAG system that combined crawling, document processing, embeddings, FAISS retrieval, a FastAPI endpoint, and a Next.js interface so users could trace answers back to original campus resources. I also built a diff-grounded LLM research pipeline that generated pull request descriptions from repository evidence and evaluated the output against raw code-change context. Those projects taught me that useful AI systems are not only about calling a model; they need clean data flows, retrieval, evaluation, APIs, and user-facing behavior that makes the output reviewable.

I would be glad to bring that practical, research-aware engineering style to RTI: build carefully, keep AI outputs grounded, document the system clearly, and learn from experienced engineers while contributing production software.

\vspace{8pt}

Best regards,\\
Aryan Miriyala

\end{document}
